\section{Protocol}\label{proto}
The \texttt{Blossom} protocol is a comprehensive system including algorithms that manage the receiving, sharing, and finalization of network updates (transactions), as orchestrated and managed by a decentralized body of trustless members. Any update made to the network is guaranteed as valid by a supermajority of all network members, with said members storing a local replica of the global network state, with the expectation that the state is cryptographically signed by a supermajority of all nodes. With this embedded decentralization, the system is highly redundant requiring greater than \(\frac{n}{3}\) concurrent faults to occur for the system to pause validation.

The procedures we have embedded within this system are intended to provide cooperative performance, with each additional node providing greater value to the system rather than accounting as losses potential, while the error tolerances of the protocol only allow for guaranteed positive behavior with the expected occurrence of false-negatives (type II errors). This strategy allows us to reduce the size and complexity of the code base by removing redundancies which guarantees admission of no false positives (type I errors).
As previously mentioned, similar to other protocols which are expected to operate in a trustless environment, each validator is required to store a replica of the network state. Each validator then references this local state during transaction verification and consensus. Due to this framework, all behavior that we will describe in this section is understood as replicated across all validator instances.

\subsection{Structure}\label{struc}
\texttt{Blossom} can be implemented on machines across the globe, with variability in network bandwidth and node compute potential, thus there is a need to optimize the protocol to mitigate issues associated with said variability. Our optimizations minimize network traffic during data dissemination while allowing for parallelization during transaction verification. The goal of these improvements is to reduce transaction time-to-finality (latency) and improve transaction processing efficiency, thus increasing the potential network throughput (TPS).

The protocol structure, resultant of these technical improvements, incorporates a  system of three (3) distinct serialized operations: \(\langle 1 \rangle\) block formation, \(\langle 2 \rangle\) block propagation, and \(\langle 3 \rangle\) transaction validation. These processes can be understood as serialized stages of the protocol, requiring the previous stage to be completed for the next to begin. While tasks within each stage can be optimized with parallelization, these discrete phases have no such parallel potential. The characteristics of each stage can best be understood as the following:

\begin{description}
\item[Block Formation] With \texttt{Blossom} every network member is able to receive transactions and form blocks in parallel. This is possible given the ‘leaderless’ nature of the protocol, with the system designed to allow independent members to reach a shared state. When a transaction is received by a member the transaction is \textit{pre-validated}. \textit{Pre-Validation} requires the member to check the transaction against the local state, as to return a preliminary verification of the transaction. Verified transactions are then added to the member's local block, with the block closed once a predetermined cap has been reached (1 MB).

\item[Block Propagation] As every network member is in possession of a unique block of transactions, this stage of the protocol deals with the dissemination of said block in the most efficient manner. Our solution involves a recursive process of sub-quorum consensus, with each member participating in a quorum (q) of size \(q = 6\). We have selected six (6) for our quorum sizes as it is the largest value which allows for both a supermajority (s) to be found (\(s=\frac{2}{3}\)) while presenting a fault-tolerance (f) greater than one (1). Each quorum undergoes a procedure of message passing, with each message furthering consensus among all quorum members. The distinct message types are: \textit{dispatch}, \textit{echo}, \textit{verification}, \textit{proposal}, and \textit{commit}. Once a quorum is completed — each member has received the unique block(s) from their peers and agreed on a shared outcome — the member moves into the next quorum where the process is repeated. The membership and order of each quorum is deterministically computed before block propagation begins, as to guarantee a shared outcome. Once a member has iterated through every quorum, they should now possess the collective data of every member on the system.

\item[Transaction Validation] To process the newly collected transactions in a concurrent manner we implement an ordered data structure (B+ tree) for the storage of blocks, with transactions stored in a \textit{B+ tree} within each block. Given the serialized nature of transaction validation, we have worked to implement some parallelization within this procedure. Acknowledging that signature verification is the process with the greatest cost (\(\sim 100,000\) compute cycles per signature), we looked to implement parallelization for this computation. Given the serialized process of updating the ledger is less costly per transaction (\(\sim 100\) compute cycles per signature), every parallel thread verifying signature correctness drastically reduces \textit{transaction validation} latency.

\end{description}

The cumulative outcome resulting from successfully processing these three (3) stages is representative of a full protocol cycle, with the outcome of said cycle validating the state unification of a decentralized and trustless body. The colloquial terminology which numerically represents this protocol cycle is understood as the protocol \textit{Block Time} (\(B_t\)).\footnote{\textit{Block Time} is the measure of time it takes for network members to create, share, and validate their block amongst the global network.}
The stages described above will be discussed in further detail in subsequent sub-sections in this document with additional nuance and descriptions [\S\ref{struc:comm_model}, \S\ref{struc:con_tree}, \S\ref{struc:ledger} respectively].

\subsection{Initialization}\label{struc:init}
Before the blockchain is an active and autonomous system it must first be initialized, and in the case of a network, outage reinitialized. The procedure for primary initialization, colloquially acknowledged as “genesis,” is generally undertaken by a platform's founding members. We do not intend to stray from this norm, defining a predetermined “genesis state” with platform insiders representing the initial network members. Depending on the context of the protocol's implementation additional members may or may not be added. 

To guarantee the immediate operation of the network following \textit{genesis} the platform is expected to maintain a network size greater than or equal to the minimum quorum size of the protocol implementation (\(n \ge q\)). While for the purposes of this paper, we have implemented a quorum size of six (6) (\(q=6\)), others may adopt differing quorum characteristics. However, despite any alteration in quorum specifications, the prior rules will still apply. Once a consensus has been reached among inceptive nodes and the genesis state has been finalized, the recursive state validation of the protocol shall continue indefinitely. This operation of \textit{initialization} is expected to occur once, at the inception (genesis) of the network implementing this protocol.

\textit{Reinitialization} is implemented in the case that the protocol is forced to ‘stall’ due to the failure of nodes to reach a shared state. A ‘stalled’ network is one where, while some sub-quorums may successfully process blocks, epochs fail to be finalized by the necessary number of signatures to be considered ‘globally validated.’ The procedure for network \textit{reinitialization} is different from \textit{initialization} and requires the following assumptions to be true: 
\begin{itemize}
    \item There exist a viable number of active nodes, with the expectation that the proportion of active nodes to Byzantine nodes is within a reasonable margin.
    \item All nodes store a replica of the global ledger, which dates back a minimum number of epochs.
    \item At some point in this shared record there was a state which can be agreed upon by a supermajority of nodes.
\end{itemize}

With these assumptions in place, \textit{reinitialization} allows participating nodes to undergo a new process of consensus, where some supermajority of nodes can agree to a shared state by  adopting a previously validated state, as to restart the regular operation of the network. In the case the need arises for \textit{reinitialization} we expect a supermajority consensus to require only a minimal degree of transposition. Reducing transposition is important as for every epoch in the past the network traverses, and the data contained within said epoch(s) will be invalidated.


\subsection{Communication Model}\label{struc:comm_model}
The topology required to facilitate \texttt{Blossom} is maintained under the purview of a fully decentralized network architecture, where every node in the system is able to communicate with any other node in the system without obstruction. However, during the act of consensus, as to limit any obtrusive and sub-optimal transmissions, the network arranges itself into a highly structured distributed network. As to restate, this ‘structured’ composition of the network is a temporal arrangement exclusive to the operation of consensus, with all non-consensus behaviors undergone as direct peer-to-peer (\texttt{P2P}) communication.

The clear and delineated distinction between structured and unstructured network communication is important to understand, as both are utilized for specialized operations. \textit{Unstructured communication} refers to non-consensus exchanges between two nodes, such as data queries, transaction requests, etc. The goal of such \textit{unstructured communication} is to prioritize accessibility while reducing network latency on jobs that do not require additional validation throughout the network. \textit{Structured communication} relates to the transmission of data relevant to consensus and is implemented to optimize data transmission to maximize network throughput. This is achieved by limiting duplicate data transmissions and limiting the number of communications required to reach consensus while allotting standard tolerances across replica nodes. For the purposes of this paper, we will be limiting all future discussions regarding the behavior of communication to the purview of \textit{structured communication}.

\subsection{Consensus Tree}\label{struc:con_tree}
To facilitate the structured communication of data across the network, \texttt{Blossom} adopts algorithms that compose the existing network mesh into a deterministic pattern, as to guarantee all network members are able to independently reach a uniform topology (to be referred to as the ‘consensus tree’). As mentioned in the previous section, this network structure is designed as a highly optimized scheme with the intention of reducing latency, mitigating byzantine faults, and allowing for exact consistency across replicas. While certain concurrent behaviors, such as state replication, are required for the successful implementation of the \textit{consensus tree}, we can assume nodes successfully operate as independent replicas; otherwise, the employ of a consensus protocol would be trivial. Our description of the network topology as a consensus tree is didactic, as when the network is organized in the deterministic pattern previously referenced, the distribution of network nodes and their relationships is visually representative of a tree (see \textit{Figure \ref{fig:tree_map}}).

\begin{figure}[h!]
    \centering
    
\includegraphics[scale = 0.2]{treemap.png}
    \caption{The consensus tree visualized, as the parallel quorums which occur for a singular node to reach consensus. This diagram presents a network size of 216 (\(n=216\)), and quorum size of 6 nodes (\(q = 6\)).}
    \label{fig:tree_map}
\end{figure}

The consensus tree structure represented in \textit{Figure \ref{fig:tree_map}} is made up of an exterior perimeter, implied by the outmost branches of the ‘tree,’ with the number of branches that meet the perimeter equivalent to the total number of nodes on there network (\(n=216\)). Within the tree there then exist moments of union where many branches merge. These moments of intersection represent the occurrence of a quorum, where \textit{q} nodes come together to achieve a unified state. In \textit{Figure \ref{fig:tree_map}} there are three (3) distinct union hierarchies, a value which can be derived from the logarithm of all network members \(\lceil \log_qn \rceil \). The centroid of the diagram, represented as the moment where all independent branches reach a shared union, is representative of all network members achieving consensus.

To further elaborate on the mechanics of the consensus tree: the data structure is a deterministic algorithm that allows each node to concurrently determine the number of quorums required to reach a shared state (i.e. consensus), as well as determine the specific nodes within each quorum. Given some arbitrary number of nodes (\textit{n}) and assuming a quorum size of six (6) (\(q=6\)), we can estimate the number of quorums necessary to reach consensus as \(\log_qn\). We subdivide the network into quorums, with each node participating in a minimum of one (1) quorum, as a means of simulating global P2P communication without the need for true P2P communication. In assuming correct behavior across all nodes, in a classical P2P system, the number of messages required to achieve consensus is equivalent to \({m(n^2-n)}\), where the minimum number of messages sent by each node to achieve consensus is represented by \textit{m}. Given an increase in network size induces an exponential increase in network load, these protocols are unable to scale, and thus untenable for a decentralized network — such as the issue for PBFT and other similar BFT-based protocols. In subdividing a network into quorums of limited size, such as those of six (6) nodes, our message overhead is reduced to \(m(q^2-q) \times \lceil \log_qn \rceil\). Restructuring the network in this manner the communication cost for each node is reduced from \(\mathcal{O}(n^2)\) to \(\mathcal{O}(\log n)\). The efficiency of this communication design, in combination with our parallelized network model, allows us to offer immense gains in network throughput when compared to previous systems.

% \begin{figure}[h!]
    \centering
    \begin{tikzpicture}
        \begin{axis}[
            scale only axis,
            axis lines = left,
            % ymode=log,
            xlabel = $n$,
            ylabel = {$\Sigma m$},
        ]
            % %Below the red plot is defined
            \addplot [
                domain=0:100, 
                samples=100, 
                color=red,
            ]
            {exp(x)-x};
            \addlegendentry{\(\mathcal{O}(n^2)\)}

            %Below the blue plot is defined
            \addplot [
                domain=0:100, 
                samples=100, 
                color=blue,
            ]
            {x};
            \addlegendentry{\(\mathcal{O}(n)\)}

            %Below the blue plot is defined
            \addplot [
                domain=0:100, 
                samples=100, 
                color=black,
            ]
            {log2(x)};
            \addlegendentry{\(\mathcal{O}(\log n)\)}
        \end{axis}
    \end{tikzpicture}
    % \caption{The rate of change in message complexity between \texttt{P2P} (\(\mathcal{O}(n^2))\), improved \texttt{P2P} (\mathcal{O}(n)), and \texttt{Blossom} \mathcal{O}(\log n)}
    \label{fig:message_rate}
\end{figure}

\begin{figure}[h!]
    \centering
    \includegraphics[scale = 0.07]{consensus_diagram.png}
    \caption{The quorum distribution scheme within each epoch, assuming a quorum size of \(q=6\) and network size of \(n=36\).}
    \label{fig:consensus_diagram}
\end{figure}

The mechanics which allow us to provide the aforementioned efficiencies are best exemplified in \textit{Figure \ref{fig:consensus_diagram}}, which displays two matrices of size \(n=36\) composed along six (6) rows and six (6) columns. This equilateral distribution is representative of a network model, given a quorum size of six (6) and a total network size of greater than two (2) quorums \(n \ge 12\). The diagrammatic color gradient is meant to represent the quorum to which each node is party (“Quorum 1”, “Quorum 2”, “Quorum 3”, etc…). Assuming each node enters the first round of consensus (as seen on the left matrix) with some unique data, they would then distribute said data with their quorum peers, with their peers equally required to share any unique data they possess. Assuming all nodes behave as expected, the members of each quorum should reach a shared state. It can then be assumed that each member of the quorum is an equal arbiter of said quorum's state, which can be further validated by the use of cryptographic signatures. Each node can then be entrusted to concurrently distribute their newly achieved state with other nodes. This is when we move on to the second round of consensus (as seen on the right matrix) where the members of each quorum have been changed. Unlike the first round where quorum members were each units in the same row of the matrix, the second round has quorum members derived from the same column. This transposition is important, as we can guarantee that every node in the first round is unique to every node in the second round (i.e. there exists no redundancy in node selection). If there is redundancy in the system, it would be due to a non-optimal number of network nodes (\(n \neq 6^{\mathbb{Z}}\)). Given all nodes are behaving as expected, each member of the second round would provide their quorum with the shared outcome of their previous round, with the end result achieving a newly shared state and thus a global consensus.

The details of data dissemination will be discussed in full in a future section (\textit{block propagation} [\S\ref{struc:block_form}]); however, for sake of clarity the following description will summarize the scheme: When a node enters a quorum they may or may not be in possession of some unique dataset (\textit{S}) which may be relevant to their peers. In the case they do possess such data, they will share their dataset in full with their peers and said peers will do the same. Each member will then merge their received datasets via a union operation as to achieve a shared outcome. Given a matrix \(X=(x_{ij})_{1\le i,j\le 6}\) and some quorum (\(q_i\)) on the \textit{i}-th row during round one (1), and some quorum (\(q_j\)) on the \textit{j}-th column during round two (2), the operation is represented as the following:
\[
q_j = \bigcup^6_{m=1}
\big (
S^j_m:\bigcup^6_{m=1}S^i_m 
\big )
\]

As mentioned previously, the outcome of the above operation should see all participating nodes reach a shared state, the equality of which can be validated retrospectively. To verify concurrent data accumulation across members, each node will maintain a local \texttt{btreemap}, guaranteeing a consistent ordering of data across nodes while efficiently removing any duplicate data entries.

While the above method is applicable in a non-adversarial (i.e. trusted) environment, given the need for this protocol to function in an adversarial (i.e. trustless) environment some alterations must be implemented to guarantee correctness across machines. In our model, we assume a fault tolerance of \(f = \frac{2}{3}\). In a quorum size of six (6) nodes, this would allow for no more than two (2) Byzantine nodes to be present, or else the quorum will fail to reach a consensus. In the case more than two (2) Byzantine nodes are present in a quorum, the non-Byzantine nodes fail to convene and are required to ignore any data they received during that quorum. In such a scenario the non-Byzantine nodes are reliant on their future quorum(s) to be non-faulty to reach a consensus. Given \textit{Blossom’s} implementation in a trustless environment, to guarantee all nodes receive all messages and reach an equivalent state, we will repeat the previously described consensus operation twice. The first recursion of the procedure will allow nodes the opportunity to initialize a communication channel with their peers and attempt to reach a consensus within each quorum; however, to mitigate the presence of any Byzantine node(s) which may adversely impact consensus, a second recursion will allow for any non-optimal communication from the first round to be remediated. An example series that mitigates the presence of any Byzantine nodes would be as follows:
\[
q_{i_1}
\Rightarrow
q_{j_1}(q_{i_1})
\Rightarrow
q_{i_2}(q_{j_1})
\Rightarrow
q_{j_2}(q_{i_2})
\]

Given the presence of a reasonable number of Byzantine nodes (\(\le \frac{n}{3} \)) by implementing the above series a supermajority consensus (\(\ge \frac{2n}{3} \)) can still be achieved, while any changes made to the network during said consensus can be finalized deterministically at the end of the present epoch, or by the completion of the following epoch (these statements are validated in the \nameref{proj} [\S\ref{proj}] section of this document).

Referencing \textit{Figure \ref{fig:consensus_diagram}}, an example membership distribution from quorums \(q_i\) and \(q_j\) from the perspective of node [0,0] would be as follows:
\[
q_i = \big[
[0,0],[0,1],[0,2],[0,3],[0,4],[0,5]
\big]
\]
\[
q_j = \big[
[0,0],[1,0],[2,0],[3,0],[4,0],[5,0]
\big]
\]

The novel algorithms we implement to yield a useful \textit{consensus tree} include our \textit{member shuffling algorithm} (MSA) and quorum selection algorithm (QSA), with the compute cost of said algorithms being \(\mathcal{O}(n)\) and \(\mathcal{O}(\log n)\) respectively, where \textit{n} is representative of the network size (i.e. number of network nodes). The MSA, while important for \texttt{Blossom} implementations on public networks with non-restrictive communication, is a relatively expensive procedure inefficient for protocol implementations on private networks.

To construct the \textit{consensus tree}, the QSA derives a non-biased assembly with the ability for said assembly to be ‘biased’ (i.e. randomized) by implementing the MSA. To efficiently randomize the data structure, the MSA implements the \textit{Durstenfeld Shuffle Algorithm} \cite{Durstenfeld1964}. This randomization is made deterministic by the implementation of a cryptographic seed, with the seed in question queried from the hash of the previous epoch — the process of computing the \textit{consensus tree} is expected to occur once each epoch, thus the use of the previous epoch for seeding this computation is both highly secure (the epoch hash derives \(2^{256}\) unique variations) and serves as a secondary concurrency guarantee (nodes with variable states should construct differing \textit{consensus tree’s} thus appear faulty to their non-faulty quorum members).

\subsection{The Ledger}\label{struc:ledger}
The ‘ledger’ is a globally consistent record of platform-specific data that facilitates the accounting of transactions and tracking of network assets in a distributed manner. Every node on the network is expected to maintain a replica of the global ledger, with updates validated by a supermajority of network participants — with said endorsements guaranteed by the use of cryptographic signatures. The ledger is structured in the form of an append-only sequential database, colloquially known as a “blockchain,” with the goal of ordering write requests atomically without collisions. To properly manage sequential write requests the ledger is designed as a NoSQL database, managing memory by implementing a data structure similar to an LSM Tree \cite{Strauch2011-hy}\cite{Oneil1996-kb}. 

The novelty of our ledger implementation is a result of our unique consensus protocol, allowing for blocks to be formed in parallel. To guarantee a consistent atomic state, despite the presence of parallel blocks, ledger additions are organized by major and minor data structures. Minor data structures refer to individual blocks which store some number of transactions, while major data structures refer to individual epochs which store some uncapped number of blocks (loosely associated with the number of participating nodes). The goal of this structure is to serialize transactions, as to be concurrently processed by independent actors.

For our specific ledger implementation the term “blockchain” is imprecise as the blocks we manage, while sequentially ordered, are independent of their fellow blocks. To elaborate  on our unique architecture from preceding systems we can look to the Bitcoin protocol; within the Bitcoin protocol, every consecutive block has embedded within its header the hash of the previous block (see \textit{Figure \ref{fig:bitcoin_tx}}) \cite{Nakamoto2008}. For \texttt{Blossom} every block embeds some sequential data within its header; however, rather than embedding the hash of the previous block we embed the hash of the previous epoch. This difference is important as the order of blocks within each epoch is reliant on the alphabetization of said block’s hash, rather than a time-stamping mechanism. To verify transaction order, block order, and epoch order, we use novel mechanisms which will be discussed in the following sections (see \textit{\nameref{struc:block_form}} [\S\ref{struc:block_form}], \textit{\nameref{struc:block_prop}} [\S\ref{struc:block_prop}], and \textit{\nameref{struc:tx_valid}} [\S\ref{struc:tx_valid}]).

\begin{figure}[h!]
    \centering
    \includegraphics[scale = 0.325]{bitcoin_tx.png}
    \caption{Bitcoin’s method for adding a timestamp to blocks, by inserting the hash of the previous block within the header of the consecutive block \cite{Nakamoto2008}.}
    \label{fig:bitcoin_tx}
\end{figure}

\subsection{Block Formation}\label{struc:block_form}
When an outside party (the “client”) sends a transaction to the network, said transaction is routed to one of the many validator nodes on the network, with the expectation that the validator will distribute the transaction to their peers. Prior to the distribution of the transaction, it is \textit{pre-validated} by the node; a process which allows preliminarily verify the validity of a transaction so as to provide an immediate response to the client. This response is non-binding as the transaction in question may not be included in the final ledger; however, within a reasonable margin of certainty, the client can move forward with the expectation that the transaction will be finalized as valid by the network.

To provide a preliminary transaction verification, the node in question who receives the transaction references their local ledger, and given the contents of the transactions hold up to scrutiny, the transaction is \textit{pre-validated}. The qualities of the transaction which are scrutinized during this process include verification that the public keys of both the transaction sender and transaction receiver are present in the ledger, that both individuals maintain the proper quantity and or class of assets required to facilitate the transaction, the epoch nonce of the transaction is correct and the cryptographic signatures which validate the individuals as signatories of the transactions are valid.

Following approval, the transaction is then added to a LinkedHashMap (the ‘transaction map’ or ‘LHM') data structure (e.g. data structure which maintains insertion order) maintained within the nodes local \textit{build block} \cite{oracle-si}. The build block is representative of the mutable state between block initialization and closure, where transactions may still be added to said block. Each block maintains a max size, equivalent to \(\sim 1,000\) transactions or \(\sim 200 MB\), and once this threshold is reached the block is closed. To verify the closure of a block, the block’s contents (e.g. transactions) are hashed, with the hash signed by the validator. This new hash (the 'block hash’), signature, and the node’s public key are then stored in the header of the block. To compute the ‘block hash’ \texttt{Blossom} implements are procedure similar to that outlined in the Bitcoin whitepaper (see \textit{Figure \ref{fig:bitcoin_tx}}) \cite{Nakamoto2008}. To summarize, each transaction in the block computes its unique hash, the result of which is concatenated with the hash of the previous transaction. This process is repeated sequentially across the transaction map, with the resulting hash representative of the entire block. The ‘block hash’ as it is both a unique identifier representing each discrete dataset (block, and or epoch) as well as a process that is easily repeated by auditors. To verify the legitimacy of any dataset every transaction, block, and epoch is expected to maintain a cryptographic signature. Given the hash (m), public key (\textit{P}), and signature (\textit{S}) are all present in the header of the said dataset, any actor can validate the legitimacy of the stored information.

\begin{figure}[h!]
    \centering
    \includegraphics[scale = 0.28]{tx_flow.png}
    \caption{The procedure of generating transaction hashes (‘Tx..’) from the body content of each transaction (‘Tx.bd’), as to be concatenated with a previous hash (‘+’) to generate a block hash.}
    \label{fig:tx_flow}
\end{figure}

\textit{Figure \ref{fig:tx_flow}} visualizes the previously described process of cumulatively concatenating transaction hashes, as to generate an overarching block hash. We considered this process to be generalized, as it is able to tokenize and sequence any possible transaction type while guaranteeing the authenticity of each transaction using embedded cryptographic signatures. If additional specificity is required by auditors, such as revealing the unique hash generated by the validator at each transaction instance, a separate sequential data structure, such as a vector, may be stored within each block. This new data structure would maintain the same sequencing as the original \textit{transaction map} with the difference that at each index position rather than displaying the unique transaction hash, the data structure would store the concatenated transaction hash. This vector could then be signed separately or with the other block data. 
Once signed a block is perceived as closed and ready to be distributed across the network. This is valuable from a network perspective as this closure mechanism asynchronously signifies the moment at which a block is ready to be distributed. Without some central arbiter or timing mechanism it would be difficult to classify the sequence of block distribution, and by handling this behavior asynchronously, we can distribute blocks in a more efficient manner. However, this true asynchronous behavior is only maintained from the initial distribution of a nodes block, as some degree of synchronicity is required during \textit{block prorogation}, which will be discussed in detail in the following section [\S\ref{struc:block_prop}]. As to summarize, dependent on the outcome of the \textit{consensus tree}, the node will be provided a list of quorum peers. The node will then distribute the blocks they manage locally with their quorums, with the goal that all non-byzantine nodes accumulate the collection of blocks formed in parallel across the network.

\subsection{Block Propagation}\label{struc:block_prop}
\textit{Block propagation} is the method by which a \textit{consensus tree} is implemented to efficiently distribute independent blocks amongst network members, to quickly reach a shared state (i.e. consensus). As per the mechanics of the \textit{consensus tree} a subset of nodes (i.e a quorum) are selected, the members of whom communicate to reach a shared state. This communication is undertaken with each node dispersing standardized messages in a serialized order, the fulfillment of which verifies the completion of their communication. The messages in question, in the order they are expected to occur, are: \textit{Dispatch}, \textit{Echo}, \textit{Verification}, \textit{Proposal}, and \textit{Commit}. These message types each play an important role in guaranteeing consistency across independent machines, while also introducing a degree of redundancy to mitigate faults. The characteristics of each message are expected to be derived from the content of the previous message(s), with this trait in line with our assumptions regarding \textit{state determinism}.

When a message is distributed, it is the responsibility of each receiving actor to verify the content of the message. This is a consequence of the protocol's decentralized nature, with any node able to communicate with any other node. However, given the computational overhead of each node increases linearly with the size of the network \(\langle \mathcal{O}(n)\rangle\), we must limit the compute spend on incorrect or invalid transactions. To optimize the protocol, each message implements a header-body format, with universal message content stored in the header and message unique content stored in the body. By standardizing this format across message types, as well as requiring each header to store: the senders public key, the last epoch hash, the current epoch nonce, the current height (also known as the ‘round’) of the consensus tree, and the message signature, if any of these data types are found to be inaccurate when a message is received, the message will be ignored with needing to process the body of the said message.

When message headers are verified affirmatively, the message will then be managed as is deemed appropriate given its typology. To discuss each message in the context of their independent qualities, as well as their role within the greater network, they will be described as follows:

\begin{description}
    \item[Dispatch] A \textit{dispatch} is a means by which each actor may distribute any accumulated ledger amendments (e.g. block(s)) with their quorum peers. This block data will be shared with the necessary signatures, as to verify the message is accurate while maintaining that each \textit{dispatch} is but an appendant to previous rounds of consensus. To dispense signatures properly, while optimizing for the high compute cost of signature verification, we implement a novel data structure termed the \textit{signature tree}, reducing the time complexity of block verification from \(\mathcal{O}(n^2)\) to \(\mathcal{O}(n\log n)\).
    \item[Echo] Within the classification of an echo message there exist subtypes for given scenarios, with those types being: \textit{Echo Response}, \textit{Echo Request}, and \textit{Echo Redispatch}.
    An \textit{echo response} is a means of signaling to the quorum that a node received a \textit{dispatch} before doing any work to validate the content of the dispatch. In optimal scenarios, this is the only \textit{echo} type shared between nodes. The latter \textit{echo} types are required given non-optimal behavior, such as when a node receives an \textit{echo} response from a peer, but the node has not yet received the \textit{dispatch} message associated with the \textit{echo response}. Recognizing that they may be missing this crucial data, the node will query the \textit{echo} sender with an \textit{echo request}, to receive the original unadulterated \textit{dispatch} message from the peer. The peer is then expected to return an \textit{echo redispatch}, sharing the missing message with the relevant node. We will discuss the additional benefits of \textit{echo} messages later in this section.


    \item[Verification] A \textit{verification} shares which block(s) a node believes to be valid or invalid from those they received via \textit{dispatch}, as well as from the block(s) already in their possession. To validate the legitimacy of a block, the block header is referenced as to determine whether a block’s hash is unique, whether the block has the proper state credentials (epoch nonce and epoch hash), and whether the signatures attributed to the block are valid. The resulting outcome of this validation process is then hashed and signed with the outcome shared within the body of the \textit{verification} message.

    \item[Proposal] A proposal is the accumulated outcome of \textit{verification} messages, as perceived by each node, the result of which should display an affirmative or negating consensus. An affirmative consensus would occur in the scenario that a supermajority of quorum peers (e.g. \(\ge \frac{2q}{3}\)) share an identical \textit{verification} message, in that the block(s) determined valid and or invalid are identical. A negating consensus would occur when an affirmative consensus is impossible  to occur (e.g. \(> \frac{q}{3}\)), such as due to a supermajority disagreement. Whatever the result, a \textit{proposal} is a means of compressing the relevant \textit{verification} messages into a concise and verifiable outcome of either \texttt{true} or \texttt{false}.

    \item[Commit] Similar to the \textit{proposal}, a commit is the accumulated outcome of \textit{proposal} messages, as perceived by each node. However, unlike the \textit{proposal} message which is used to share the perceived outcome of each node, a \textit{commit} is a definite statement from each node affirming their future behavior, whether it is to operate under the expectation of a \texttt{true} or \texttt{false} consensus. \texttt{True} would signify a consensus was reached during the quorum and some collection of blocks is expected to be ratified. \texttt{False} would signify a consensus was not reached and all novel content from the quorum should be ignored.
\end{description}

The outcome of each quorum is dependent on the type of nodes present in each quorum, as well as the context within said quorum comes together. If enough Byzantine nodes are present within a quorum there will be no outcome, despite any consensus between non-Byzantine nodes, that will be committed as true. A similar non-consensus can occur if the nodes that are placed together within a quorum are less performant than the greater network, the consequences of which will have said nodes forced to prematurely disperse. These behaviors, while non-optimal for any non-Byzantine and performant nodes within each quorum, provide a means of limiting the consequences of any possible faults while allowing systems to fail gracefully in a performant manner. It is here where we must discuss the consequences of faults that occur before or during each quorum, and the methods we adopt to mitigate said faults within the protocol.

As a consequence of this, \textit{cooperative consensus} model there exist two categories of faults that are the predominant complications for this style of the protocol. The first is a consequence of the need for synchronicity (termed “synchronicity failures”), and the second is a consequence of the distributed nature of the need for block accumulation (termed “distribution failures”). To describe each failing appropriately we will detail the variations in the following sections (\textit{\nameref{struc:sync_fail}} [\S\ref{struc:sync_fail}] and \textit{\nameref{struc:dist_fail}} [\S\ref{struc:dist_fail}]).

The measures described in the above sections are meant to outline the solutions we implemented to solve the most common of possible faults which may occur during \texttt{Blossom}. These behaviors do not dissuade our belief that other types of faults may occur, and we expect that as additional testing is undergone such issues may arise, but at the same time, we expect that \texttt{Blossom's} unique scalability incentivizes a widely decentralized network of non-Byzantine nodes, the existence of which substantiates a lively and performant network, which far exceeds the present outcomes of other precedent systems.

Following the successful procedure of block propagation the node should be party to the aggregated dataset (e.g. blocks) accumulated from each distributed non-Byzantine node on the network. This new dataset (otherwise referred to as the ‘unverified state’ or ‘preliminary epoch’) will then be required to be validated in its entirety, transaction by transaction.

\subsubsection{Synchronicity Failures}\label{struc:sync_fail}
Due to the need for quorum members to reach a synchronous state to undergo cooperative consensus, it is possible that variances in node hardware, or the existence of Byzantine nodes, introduce undue consequences such as the inability for nodes to progress beyond dispatch messages. The reasons for such a scenario to occur are several: a node lacks the bandwidth necessary to receive dispatch messages in a reasonable time, a Byzantine node never shares a dispatch message, or a node is still participating as a member of a previous quorum and unable to share a dispatch message. Any of these scenarios would cause a quorum to hang at the dispatch message stage, with the fault due to the need for synchronicity. To mitigate non-performant or Byzantine nodes we implemented the echo message. As mentioned previously, this message type is a redundancy feature used to notify the quorum that a node has received a given dispatch message. This is a valuable tool as it alerts nodes when they are missing content, as well as provides a model for each quorum member to independently classify Byzantine nodes from non-Byzantine nodes. This solution can be best represented by the use of a matrix of size \(q^2\) (see \textit{Figure \ref{fig:consensus_diagram_fail}}), with each row of the matrix representing a single quorum member, and each column of the matrix representing an expected message type. Each node is expected to ignore their own row, as they can assume they are sending and receiving the proper messages. The remaining items of the matrix will be filled in when the node receives the appropriate message from the relevant quorum member. In a scenario where all nodes are behaving properly (see \textit{Figure \ref{fig:consensus_diagram_fail}} (\textit{Left})), we will see that all items of the matrix are storing a ‘True’ value. However, in the case that some Byzantine nodes exist (see \textit{Figure \ref{fig:consensus_diagram_fail}} (\textit{Right})) we can see that a pattern emerges where faulty nodes send both no dispatch messages nor echo messages (noted as ‘…’, e.g. ‘null'). By implementing a reasonable wait mechanism to insure false nodes are not rejected prematurely, we can simply isolate byzantine nodes within each quorum. In the scenario that a non-byzantine node is unable to send a dispatch message, as they are currently a member of a preceding quorum, there are several possible solutions. As to mitigate a node appearing faulty to their peers, despite lagging behind their quorum members, the node may still send an echo message despite not yet sending a dispatch message. This behavior signifies that the node is delayed in a previous quorum and is not faulty. In the case that a supermajority of nodes (\(\frac{2q}{3}\)) have already shared a dispatch message, and given the node in question is a member of a non-consensus quorum, or stuck before the verification message round of their present quorum, the node is then required to send their dispatch message within a reasonable time. In the scenario that a node is near the completion of their quorum (this is signified as the node has already transmitted their verification message to a possibly consenting quorum), the following quorum members must then wait for the lingering node, until they are then called upon by a following round. If the same scenario was to occur, but this time over the course of two epochs, then the idling quorum members would be required to wait until they are called into their subsequent quorum.
\begin{figure}[h!]
    \centering
    \includegraphics[scale = 0.3]{consensus_diagram_fail.png}
    \caption{(\textit{Left}) A messaging matrix with no observed byzantine nodes (where ‘true’ represents non-Byzantine behavior). (\textit{Right}) A messaging matrix with a byzantine node at the position of node \textit{F} (where ‘…’ represents Byzantine behavior).}
    \label{fig:consensus_diagram_fail}
\end{figure}

\subsubsection{Distribution Failures}\label{struc:dist_fail} In the case a quorum fails to reach consensus, any non-Byzantine nodes present within said quorum will have the undue consequence of missing crucial block information intended to be shared by other non-Byzantine nodes in their quorum. This will result in these non-Byzantine nodes finishing the epoch with missing data, thus reaching a different state from other members of the network. The measurement to determine the frequency of this scenario is termed the failure bleed and is calculated as the number of non-Byzantine nodes which lose state as a consequence of their association with Byzantine nodes (relative to the proportion of Byzantine nodes on the network). We previously stated our solution to this issue (see \textit{\nameref{struc:con_tree}} [\S\ref{struc:con_tree}]); however, there are additional details missing in that description. As to elaborate on our solution, to mitigate the possibility that any node may be missing data as a consequence of their quorum(s) not reaching consensus, we implement a secondary consensus tree recursion where nodes reconvene in the same quorums as they were previously members. The intended outcome of this procedure expects that by allowing the node to interact with at least one (1) quorum which previously achieved a successful consensus, the node will be able to ascertain any data they missed. While it remains possible this procedure fails to provide the node with all missing data, such a possibility is statistically insignificant (relevant to the proportion of Byzantine nodes present in the network). The part of this solution which was not stated previously was that, unlike the first recursion of the consensus tree, which implemented the consensus messages: Dispatch, Echo, Verification, Proposal, and Commit; the second recursion implements more performant message types specialized for the infrequent need to update a minority of nodes. The messages in question, in the order they are expected to occur are: Appraisal, Echo, Verification, Proposal, Request, and Commit. The outcome of this novel message pattern is a reduction in infrastructure costs (bandwidth and latency) given all nodes share an equivalent state, while efficiently facilitating the distribution of blocks (given the need arises). To manage these potentialities it is the goal of each node to share their appraisal message, which is limited to exclusively storing the hash and relevant signatures of each block in their possession. This drastically reduces the data overhead of each message, thus increasing communication efficiency. The \textit{echo} message operates identically to the first recursion of the consensus tree, functioning as a redundancy measured to guarantee all messages are received and no additional faulty nodes are present. The \textit{verification} message is identical to that in the previous recursion, with the goal of presenting some consensus among nodes. The \textit{proposal} message, as  done in the previous recursion, aggregates the verification messages to return a finalized consensus. The \textit{commit} message, identical to the previous recursion, shares the outgoing behavior of each node (true or false). The second novel message, the request message, is implemented when a non-byzantine node recognizes there is a consensus among their peers, which they are not a part of, and requests to receive the blocks they are missing. This request is shared globally as to receive partitions of the missing blocks from each node. Following the request-response a commit is shared. In the best case scenario the second recursion adds a minimal cost equivalent to the latency of five (5) messages (Appraisal, Echo, Verification, Proposal, and Commit), while in the worst scenario, we must account for the latency of the additional messages, and dissemination of requested blocks.

\subsection{Transaction Validation}\label{struc:tx_valid}
Following the completion of \textit{block propagation} each node is expected to maintain a unified dataset, a constituent of all blocks distributed by every non-byzantine node on the network. This information substantiates the \textit{preliminary epoch}, which must now be verified independently by each node as to reach a finalized state. This finalized state is represented by an ‘epoch hash’ which simply confirms the unique and newly ratified condition of the ledger, verified by each node by use of a cryptographic signature. Despite maintaining an identical dataset, the protocol requires each validator to verify their \textit{preliminary epoch} independently, as the consensus proof is not based on the work of each node, but the shared outcome of all nodes (i.e. the supermajority agreement of all independently achieved epoch states). If a method was implemented which returned an identical consensus proof, while not requiring each node to verify the entirety of each preliminary epoch, but rather a minority partition of the preliminary epoch, our protocol could achieve additional performance increases. Given our assumption regarding \textit{state determinism}, as long as all nodes are provided an identical data input, we can expect all non-byzantine nodes to concurrently reach an identical state. If some  supermajority of nodes reached such a shared state, the epoch in question would be considered provably correct, and pass as validated.

The process to reach a finalized state requires each node to verify the legitimacy of every block and every transaction within the \textit{preliminary epoch}. The process for verifying blocks and transactions is slightly different; however, the positive outcome of these procedures allows each node to affirm the state of their finalized epoch. We have discussed the behavior required to validate transactions and blocks in previous sections (referred to as \textit{pre-validation}), but given the finality of this new procedure, the outcome of this process will impact the local ledger of each node, thus requiring additional novel validation proofs.

Prior to discussing the specific behaviors we implement during verification, it is valuable to summarize the data structure of the \textit{preliminary epoch}, as it relates to the process of \textit{transaction validation}. The goal of \textit{block propagation} is to allow for the blocks formed in parallel by each node to be efficiently distributed, as to inform every other node of all network appendages. However, given the requirement for concurrent computation across decentralized actors, we guarantee computation determinism by implementing a B-tree data structure during the local storage of blocks. This allows each block to be identically sequenced across actors, assuming equivalent input data was provided following \textit{block propagation}. This uniform sequence of blocks helps induce \textit{state determinism}, given no other faults occur during the verification process. By removing any uncertainty derived from a lack of concurrency among actors, disagreement among nodes is increasingly unlikely, and we can expect byzantine nodes to occur at intervals equivalent to our assumption regarding \textit{byzantine faults}, or at other rates due to some uncertain malicious action.

The formal process we implement to reach an ‘epoch hash’ requires nodes to ‘verify’ the legitimacy of each transaction in a non-ordered manner (i.e. verifying transaction signatures, epoch nonce and hash, etc…), followed by the sequential verification of transaction requests to alter the ledger. We have chosen to isolate signature verification, as well as other credibility proofs, due to their relatively high compute cost and necessity for transaction approval. Once the ‘credibility’ information of a transaction has been verified, the content of said transaction may be verified as to update the ledger.

As to clarify, when we mentioned ‘updating the ledger’ we are referencing the updating of some key-value pointer, rather than the altering of some mutable dataset. Due to the ‘append only’ nature of the ledger, each transaction is stored in a time-sensitive manner relative to their predecessors. For this reason, the ledger is not one, but two data structures; the first is a time series data structure where each transaction points to its predecessor; the second is a simple hashmap with the value of each index updated given the outcome of each transaction. 

The process of verifying a transaction's credentials, otherwise known as the transactions header information, is referenced as occurring in a ‘non-ordered’ fashion as the verification of this information is not required to occur in some deterministic order across the entire dataset. The only information required to be processed in a distinct order is the ledger appended information, otherwise known as the transaction body, with each transaction reliant on the outcome of the prior transaction.

The specific serialization process we implement to reach a verified state required each node to iterate over their stored B-tree in sequential order. This process is efficiently completed with the use of a B+Tree, which is optimized to allow for direct transposition between leaf nodes' memory. Starting from the first leaf in the dataset, with the block hash operating as the key of the leaf, the full block associated with that key is referenced. Similar to each transaction, a block maintains a header and body, with the header containing credential information while the body maintains novel transactions. The header of each block should have been previously verified during \textit{block propagation}; however, in the case additional redundancy is required it is possible to verify the header of each block again. With a block verified as correct, the transactions stored within each block may now be verified in sequential order.

Given the block content is affirmed following reverse verification, the \textit{transaction map} may now be iterated over to check for correctness. Assuming the header information of each transaction was previously approved, it would be the responsibility of the validator to sample the body information of each transaction sequentially, as to see if the information endorsed as part of the transaction is feasible, given the present state of the ledger. A transaction may fail at this stage given the wallet constituent to the transaction does not possess the necessary tokens, or if any other information related to the transaction is incorrect. However, in the case, the transaction passes all stages of validation, the validator appends their ledger to accommodate the newly approved transaction.

It is important we rigorously follow the sequence of transactions and blocks when updating the ledger for two important reasons. The first is that some determinism is required to allow for future auditors to reverse engineer the \textit{transaction validation}, as to verify the state of the present epoch; the second is that due to the parallel nature in which blocks are formed it is possible that two independent transactions, depending on the order those transactions are read, may or may not conflict. An example of this behavior would be: in the case that an individual begins the epoch with a starting balance of 100 units, and during the epoch, they post two transactions. In the first transaction, the individual receives 50 units, and in the second transaction, the individual sends 120 units. While from the temporal perspective of the individual, they were able to send 120 units as the net outcome of the two transactions would leave them with 30 units. However, from the perspective of the network, there is a 50\% likelihood that the transaction sending 120 units will occur before the transaction receiving 50 units in the eyes of the protocol. In the case, such an order occurs the network would see the transaction sending 120 units as invalid as the user would have a negative balance in their account, and as a result, the transaction would be invalidated. Due to the likelihood of such scenarios, it is important transactions are processed in a concurrent, similar, and deterministic manner.

As each block in a given epoch was created in parallel, we have not yet discussed the methods implemented to generate an ‘epoch hash,’ which accounts for the ordering and verification of said blocks and their transactions. While one may assume we implemented the previously discussed concatenation procedure, due to the ability for invalidated transactions to be passed within a given block, we must introduce a method for handling this ambiguity. Following \textit{block propagation}, if for any reason a transaction is invalidated, the ledger will not be appended with said transactions query; however, despite failing to be approved, we do not remove the transaction from the block, nor do we remove the entire block from the epoch. Instead, to resolve this issue we implemented a novel procedure we call ‘transaction tagging.’ \textit{Transaction tagging} requires some number of additional bytes to be appended to the front of every hash. This preceding byte string is intended to be unique to each transaction type (signaling to any validator and future auditor how to handle the transaction) while displaying the verification status of the transaction.

Assuming a single bit is utilized to declare the verification status of each transaction (e.g. boolean of \texttt{true} \(\langle 1 \rangle\) or \texttt{false} \(\langle 0 \rangle\)), when implemented with a single preceding byte, \textit{transaction tagging} allows for the classification of 128 (\(2^7\)) unique transactions types. If it is in the interest of our platform to allow for a greater number of transaction variations, we could instead prepend each hash with two (2) bytes, for a total of 32,768 (\(2^15\)) unique transaction types. When observed within the historic record, as a transaction can be simply determined \texttt{true} or \texttt{false} — with \texttt{true} transactions valid following all verification steps, while \texttt{false} transactions are minimally valid during block formation but invalidated by a later procedure — we maintain the integrity of the original block, while clarifying the status of the transaction.

The specialized procedure we implement to reach an ‘epoch hash’ and ‘updated ledger state’ is similar to the procedure described in \textit{\nameref{struc:block_form}} [\S\ref{struc:block_form}], with the exceptions that invalidated transactions are ignored, and the ‘core’ hash of existing datatypes (transactions and blocks) are not updated.\footnote{We define the ‘core hash’ as the part of the hash which succeeds the \textit{transaction tag}. Specifically ignoring the ‘verification status’ of each hash.} It is important we allow for the non-core (the ‘verification status’) part of each transaction to remain mutable, as that operation is at the digression of each validating node. Given the header of all transactions and blocks stored in the \textit{preliminary epoch} have been validated, the node is then required to repeat the concatenation procedure, this time utilizing the entirety of each datatypes hash; however, instead of updating the stored hash of each datatype with the outcome of this procedure, we store this mutable ‘derivative hash’ in a separate data structure. Our reasons for not updating each datatype with a new hash and accompanying signature are due to the high computational overhead for doing so, as well as the superfluous nature of the operation. Instead, similar to the formulation of a ‘block hash', each ‘epoch hash’ is computed by passively iterating over the sequence of the epoch’s embedded components (i.e. blocks and transactions). While iterating over this sequence prior to forming a finalized ‘epoch hash’ this temporary data structure is described as the ‘derivative hash.’ However, given the iterative computation is complete the resulting ‘derivative hash’ is equivalent to the ‘epoch hash.’ This process is displayed in \textit{Figure \ref{fig:verif_flow}}.

\begin{figure}[h!]
    \centering
    \includegraphics[scale = 0.25]{verif_flow.png}
    \caption{The procedure of generating a novel epoch hash (‘New Epoch Hash’) from the hash of each sequential transaction (‘Tx…’) and block (‘Block’), as to be concatenated (‘+’) with a previous hash.}
    \label{fig:verif_flow}
\end{figure}

Following the computation of a novel epoch, if one were interested in querying a specific historic transaction maintained by the ledger, there are two possible methods that could be implemented. The first option would allow for the querying actor to reference the most recent state of some index (e.g. account) on the ledger, and then iterate across all transactions sequential transactions where an update was made to that index.\footnote{This behavior is similar to that of a Skiplist \cite{Pugh1990-yz}} The second option would allow the actor to use the location information of the transaction (i.e. the epoch hash or nonce, the block hash, and transaction hash) to query the transaction. Given each transaction maintains the nonce of each epoch, but not the hash of their epoch and block, an actor could also iterate over every transaction block.
